\section{Calibration and positive masses}\label{L:part4}

The long-circle expansion permits a leading term whose coefficient could
vanish. Positive polynomial norms will first rule this out for $Q$ and $R$.
We then extract the Laplace estimate needed both for winding loops here
and for open-strip crossings in Section~\ref{L:part6}. The physical
winding amplitude will be $J_n\sim\sqrt3/(6n)$; the rooted polygon
mass will remain uniformly bounded.

\subsection{Positive norms and nonzero leading terms}
Denote homogeneous \(Q,R,Y\) by subscripts \(n\) at width \(n\). Use \(c=3/4\) and the Fourier weights
\[
 \mu(\nu)=\frac{\nu\cosh(\pi\nu/4)}{\sinh(\pi\nu)},\qquad
 w_b(\nu)=\frac{k}{\sqrt2}\,\frac{\sinh(\pi\nu/4)}{\sinh(\pi\nu)} .
\]
Indeed \(M(e^s,1)=\int e^{i\nu s}w_b(\nu)d\nu\), and \(A(e^s,1)= -i\,\mathrm{PV}\int e^{i\nu s}\mu(\nu)d\nu/\nu\), \(A(x,y)=(x^2-y^2)/h(x,y)\). The transforms follow, for example, by summing the residues at \(\nu=i j\) for \(s>0\), including the half term at 0 in the principal value (use horizontal tops at half-integer heights, then symmetry/analyticity). Write \(h_j(b')\) for squared monic polynomial norms, after orthogonalization in a weight \(b'\). Confluence in log coordinates gives
\[
 \frac{Q_{n+2}}{Q_n}=\frac{k^{2n+1}}{n!(n+1)!} h_n(\mu),\qquad
 R_n=k^{n(n-1)}\prod_{j=0}^{n-1}\frac{h_j(w_b)}{j!^2}.
\]
For the first identity take derivatives in the Pfaffian, ordered starting at degree 0. Opposite-parity entries between degrees \(a',b'\) are \(-i\cdot i^{a'}(-i)^{b'}\int\nu^{a'+b'-1}\mu\), and same parity gives zero. In even size the resulting determinant block (evens vs odds) is the Gram of successive even monomials up to row signs; in odd size the augment first eliminates degree zero, leaving the Gram of successive odd monomials. The alternating-order Pfaffian is the determinant in odd/even positions; the entry signs give overall \((-1)^{\lfloor n/2\rfloor}\), canceled by the reversed Vandermonde in \(F\). This proves the formulas and positivity.

Compare with
\(\mu_0(\nu)=\nu/[2\sinh(c\pi\nu)]\le\mu(\nu)\),
\(w_0(\nu)=k/[2\sqrt2\cosh(c\pi\nu)]\asymp w_b(\nu)\). Their Laplace transforms are respectively \((2c)^{-2}\sec^2(s/(2c))\), \(k/(2c\sqrt2)\sec(s/(2c))\), by the same residue calculation. Thus
\(h_n(\mu_0)=n!(n+1)!/(2c)^{2n+2}\). Indeed expand the two-variable Gram kernel at \(s+s'\) by the addition formula: for \(\sec^2\) it uses features \(\sqrt{n+1}\sec^2(\tau)\tan^n(\tau)\), \(\tau=s/(2c)\), times the square-root mass. Their leading Taylor coefficients give the squared monic norms by triangularity.

From \(O_Z(r)\) in the finite calculation, the unprobed raw sums at the two twists are \(R_n^2/(2k)^{n(n-1)/2}\) and \(\kappa^{2n}Q_n^4/(2k)^{n(n-1)/2}\). Hence the long-circle normalizations of \(Q_n,R_n\) themselves correspond respectively to
\[
 B_Q(n)n^{5/48},\quad B_R(n)n^{-7/24},\qquad
 B_Q(n)=(k/(2c))^{n^2/2}k^{-n/2},\quad B_R(n)=(k/(2c))^{n^2}2^{-n/2}.
\]
\begin{proposition}[Positive homogeneous normalization]\label{L:prop:positive-normalization}
There are constants $c_Q,c_R>0$ such that
\[
 Q_n\sim c_Q B_Q(n)n^{5/48},\qquad
 R_n\sim c_R B_R(n)n^{-7/24}.
\]
After these leading factors are removed, both sequences have expansions
in terms $n^{-\gamma}$ with rational $\gamma\ge0$ and polynomial
coefficients in $\log n$, to arbitrary inverse-power accuracy. The positive
constant leading terms and all coefficient polynomials are common to all
integer sizes. Any fixed number of finite differences may be taken by
first choosing the expansion to sufficiently high accuracy.
\end{proposition}
\begin{proof}
The comparison \(\mu\ge\mu_0\) (minimize the squared norm over monic polynomials) gives \(Q_n\gtrsim B_Q(n)\) by multiplying the norm lower bounds (in both parities). Thus its leading coefficient is nonzero by the gap of at least 1 (greater than \(5/12\)) in the expansion of its normalized fourth power.

For the other coefficient it is convenient to use
\[
 Y(Z)=(-1)^n\left(1-k\,(1/z)_{z\in Z}^{\mathsf T}M(Z,Z)^{-1}(z)_{z\in Z}\right).
\]
Indeed the bracket is the first cross term of the Schur complement of \(M\) between an external variable near infinity and one near zero, relative to \(M\) there. Thus it gets multiplier \(-1\) relative to the fused value and \(+1\) at a pole pair by the Schur identity in the finite proof. It has denominator cleared by \(R\) by the same determinant argument and is bounded at single-variable extremes, e.g. remove that site by one Schur elimination using \(M(x,y)\sim k y/x\) as \(x\to\infty\). The pair specializations therefore establish equality (degree \(\le2(n-1)\) after clearing), starting with sizes 0 and 1 directly. Let \(p_j\) here denote orthonormal polynomials of \(w_b\), positive leading coefficients. Then at confluence the cross evaluation uses \(\nu=i\) in both adjoint arguments, so
\[
 Y_n+Y_{n+1}=k|p_n(i)|^2 .
\]
The sum \(\sum_{j<n}|p_j(i)|^2\) is comparable to \(n^{3/2}\). It suffices by the variational evaluation bound and two-sided weight comparison to do this for \(w_0\); its kernel expansion just as above for \(\sec\) gives the orthonormal polynomials there as coefficients of
\[
 h_0(w_0)^{-1/2} (1+z^2)^{-1/2}e^{2c\nu\arctan z}.
\]
At \(\nu=i\) these are proportional to \(i^j\) times coefficients of \((1+z)^{1/4}(1-z)^{-5/4}\), asymptotic to a positive constant times \(j^{1/4}\) by the binomial theorem (the first factor has absolutely summable coefficients and their sum is nonzero; the second coefficients grow as constant times the indicated power by their consecutive ratio, with bounded backward shift ratios).

In particular \(0<Y_n\le C(1+n)^{3/2}\), and \(R_n\) cannot be superpolynomially small relative to \(B_R(n)\), using \(Y_n=\kappa^n Q_n^2/R_n,\ \kappa^n B_Q(n)^2=B_R(n)\). Its expansion thus has a first nonzero coefficient common to all sizes, and so \(Y_n\) has at least a leading power times log power description with nonzero coefficient. Summing \(Y_j+Y_{j+1}\) now forces \(Y_n\asymp n^{1/2}\) (if the normalized \(R_n^2\) first occurred at positive order \(p\), \(Y_n\) would instead grow as \(n^{1/2+p/2}\) up to logs), proving nonvanishing of the actual leading term of \(R_n\) as well. Taking the positive roots of the normalized fourth-power and squared
expansions now gives the stated expansions for $Q_n,R_n$. For finite
differences of any fixed order, first take the remainder in
Proposition~\ref{L:prop:long-circle} to an arbitrarily higher fixed accuracy;
the finite difference of that remainder stays within the required error.
This proves the proposition.
\end{proof}

\subsection{A polynomial Laplace estimate}
Let $p_j$ be the positive-leading orthonormal polynomials for $w_b$, write
$h_j=h_j(w_b)$ for its squared monic norms, and set
\[
 f_j(\theta)=\int p_j(\nu)e^{\theta\nu}w_b(\nu)\,d\nu.
\]
The following estimate retains its first correction because the leading
terms in the winding calculation will cancel after taking real parts.

\begin{lemma}[Polynomial Laplace estimate]\label{L:lem:polynomial-laplace}
The positive Jacobi off-diagonals for \(w_b\) satisfy
\[
 \sqrt{h_n/h_{n-1}}= n(2c)^{-1} s_n,\qquad s_n^2=1+\frac{d_0}{n^2}+O(n^{-2-\epsilon})
\]
for some real \(d_0\) and some \(\epsilon>0\); the on-diagonals vanish by symmetry. Put \(S_n=\prod_{j=1}^n s_j\); then \(h_0 S_\infty^2=k/(2c\sqrt2)\).
For $0<\tau<\pi/2$, the ratio below is $1+O(n^{-1})$, locally uniformly
in $\tau$, and has the pointwise asymptotic
\begin{equation}\label{L:eq:laplace-correction}
 \frac{f_n(2c\tau)}{\sqrt{h_0}\, S_n \sec\tau\tan^n\tau}
       =1+\frac{d(\tau)}{n}+o(n^{-1}).
\end{equation}
\end{lemma}
\begin{proof}
The Jacobi-coefficient expansion follows from the smooth second ratio of
$R_n$ in Proposition~\ref{L:prop:positive-normalization}; the first ratio
gives the limit for $h_0S_n^2$. The on-diagonal coefficients vanish because
$w_b$ is even. To prove the Laplace assertion, \(g_n=f_n/(\sqrt{h_0} S_n)\) at \(2c\tau\) has flow \(g_n'=n g_{n-1}+(n+1)s_{n+1}^2 g_{n+1}\) from the degree-zero unit vector. The base propagator \(U(t)\) for all \(s=1\) has entries \(\sec t\,[z^n](z+\tan t)^j/(1-z\tan t)^{j+1}\). Duhamel's equation for the ratio \(H_n\) displayed above is
\[
 H_n(\tau)=1+\int_0^\tau \tan t\
  \mathbb E_{n,\tau,t}\big[(K+1)(s_{K+1}^2-1)H_{K+1}(t)\big]\,dt ,
 \qquad
 \Pr(K=j)=\frac{U(\tau-t)_{nj}U(t)_{j0}}{U(\tau)_{n0}}.
\]
Indeed the iterated perturbation series is absolutely bounded on compact time intervals below \(\pi/2\), by positivity of base entries and bounded \((K+1)|s_{K+1}^2-1|\) (Volterra factorial bound). It equals the flow above by Taylor coefficients, which come from powers of the Jacobi recurrence. For analytic justification the base positive Taylor kernels give the same absolute majorant for the coefficient series. Here \(K\) is the sum of a binomial with parameters \(n,a/(a+b)\) and an independent negative binomial (pgf \([(1-ab)/(1-abz)]^{n+1}\)), \(a=\tan t,\ b=\tan(\tau-t)\). Hence \(n a\,\mathbb E(1/(1+K))\le a+b\). This gives \(H_n=1+O(1/n)\) locally; multiplying the integral by \(n\) gives a limit by dominated convergence and the elementary large-number limit for \(K/n\) at \(0<t<\tau\). Inverse moments used here are uniformly integrable after scaling by \(n\) at fixed \(t>0\), e.g. by the binomial exponential tail below half its mean. This proves the Laplace estimate.
\end{proof}

\subsection{Winding and rooted polygon masses}
We now apply the normalized polynomial estimates to the physical gluing
formulas. Throughout, $J_n$ counts winding polygons modulo whole-row
translations on the regular $n$-column cylinder, as in
Proposition~\ref{L:prop:periodic}.

We next evaluate \(J_n\) at the regular honeycomb point, in original conventions \(t=q^{-1}\). Confluence in the wave formula yields
\[
 J_n=\frac{(-1)^{n+1}\sqrt2}{Y_n}\Re[P_n(\pi/4)P_n(\pi/2)],\quad
 P_n(\theta)=e^{-i\theta}-\sum_{j<n}p_j(-i) f_j(\theta).
\]
This just uses projection onto polynomial degrees \(<n\) in the Fourier kernel followed by evaluation at \(-i\), with its conjugate or reflection for the other waves.

Set \(A_n=|p_n(i)|=\sqrt{(Y_n+Y_{n+1})/k}\); \(p_n(-i)=(-i)^n A_n\) by the positive-off-diagonal recurrence. For \(r=\tan(\theta/(2c))\) at each of our two angles,
\[
 P_n(\theta)=(-i)^n A_n f_n(\theta)
   \left[\frac1{1+i r}+\frac1{4n}\frac{-i r}{(1+i r)^2}+o(n^{-1})\right].
\]
For \(r<1\) this sums the tail starting at \(n\); the full series reproduces \(e^{-i\theta}\) by convergence and Taylor projection (absolute on smaller disks by the positive recurrence and estimates above). For \(r>1\) sum the negative prefix backwards, dropping the exponentially negligible free term. In both directions use \(A_{n+j}/A_n=1+j/(4n)+o(n^{-1})\) for fixed \(j\) by smooth expansion; \(f_{n+j}/(r^j f_n)=1+o(1/n)\) by the preceding estimate. The sums with first-order errors are dominated by geometric series times polynomials in \(|j|\), e.g. restrict indices first to \([n/2,2n]\) using the smooth estimates there and then exponential tails.

The two \(r\)'s are reciprocal. The product of the leading brackets is purely imaginary; the real correction is \(-2/(4n(r+r^{-1})^2)\). Using \(A_n^2/Y_n\to 2/k\) and the first ratio limit above gives
\begin{equation}\label{L:eq:calibrated-winding}
 J_n\sim \frac{2(1/4)}{c(r+r^{-1})n}=\frac{\sqrt3}{6n}.
\end{equation}

At the alignment \(t=1\), \(E_+\) is the total polygon mass through a \textbf{top} port in the physical chain by the mark reduction (both vacuum halves independent of row parameter). In original conventions the pole-pair recursion gives
\begin{equation}\label{L:eq:top-port}
 E_+=1+\kappa\left[\frac{1+q^3}{d(1,1)}\frac{Q_{n+1}Q_{n-1}}{Q_n^2}-1\right]
       = (2-\sqrt2)/3+O(n^{-2})
\end{equation}
by the smooth expansion. In particular the plane total polygon mass through a fixed port (hence also through a fixed center) is finite: any finite selection projects injectively for sufficiently large width, and we have symmetry of the edge directions.

